\section{Verification and Testing}

Ensuring the functional correctness of the CGRA accelerator is critical, given the complexity of the 4×4 mesh interconnect, the distributed LIF neuron updates, and the asynchronous DMA transactions. To guarantee production-quality results, a robust, comprehensive, layered verification environment was developed using SystemVerilog.

\subsection{Verification Methodology: Industry-Standard Rigor}
The verification strategy transitions from simple directed tests to a high-coverage Constrained Random Verification (CRV) approach. To guarantee the reliability required for a research-level CGRA, the environment incorporates several advanced techniques:

\begin{enumerate}
    \item \textbf{Protocol-Aware Assertion-Based Verification (ABV):} \\
    Utilization of \textbf{SystemVerilog Assertions (SVA)} to enforce AMBA AXI4 and APB protocol compliance. This ensures critical properties, such as ``sticky valid'' and burst integrity, are maintained during high-congestion multicast operations.
    \item \textbf{Backdoor Memory Access and Self-Healing Scoreboards:} \\
    The environment validates the Bi-Directional DMA Write Engine through direct backdoor access to the AXI RAM model. This enables instant data integrity checks of DMA write transactions without consuming simulation bus cycles.
    \item \textbf{Cycle-Accurate Golden Reference Model:} \\
    A bit-accurate SystemVerilog behavioral model that predicts saturation effect, LIF threshold crossings, and NoC routing delays with cycle-perfect fidelity.
\end{enumerate}

\subsection{Verification Flow}
The testbench follows a closed-loop execution flow:
\begin{enumerate}
    \item The \textbf{Test Case} defines a specific architectural scenario (e.g., multicast stress test).
    \item The \textbf{Sequencer} breaks the scenario into a series of abstract transactions.
    \item The \textbf{Driver} (Bus Functional Model) converts transactions into physical pin-level signal toggles according to AXI/APB protocols.
    \item The \textbf{DUT} processes the signals and drives response pins.
    \item The \textbf{Monitor} observes the output pins and reconstructs the transactions.
    \item The \textbf{Scoreboard} validates the observed transactions against the Golden Model.
\end{enumerate}

\subsection{Regression Test Suites}
The verification plan is executed through 25 specialized suites (8,915 test cases) organized into 6 CRV meta-suites:
\begin{itemize}
    \item \textbf{System Integrity:} Merges Suites A-C and R. Validates APB register access, DMA datapath, protocol compliance, and streaming with 400+ randomized vectors.
    \item \textbf{Fabric Stress:} Merges Suites L and O. Streams random data through the PE mesh interconnect, verifying pipeline integrity, parallel compute, and routing sweep with 300+ vectors.
    \item \textbf{Robustness:} Merges Suites S, W, and AA. Random fault injection, reset recovery, stall injection, and IRQ stress testing with 300+ vectors.
    \item \textbf{ISA Regression (Suite AD):} Extensive randomization of all 21 opcodes with 500 random vectors per opcode, covering boundary cases such as zero multipliers, saturation, and negative underflows.
    \item \textbf{LPR Features:} Validates RELU (opcode 19), MAX (opcode 20), hardware loop CSR programming, and east-edge result collection for the License Plate Recognition demo.
    \item \textbf{DMA Writeback (Suite AE):} Validates AXI write channel with round-trip data integrity, burst protocol compliance, 4KB boundary crossing, and concurrent read/write operations.
\end{itemize}

\subsection{Hardware-Software Co-Verification}
A critical bridge to physical deployment is the hardware-software co-verification flow. This involves executing cache-coherent C-based drivers on the dual-core ARM Cortex-A9 APU within the Zynq-7000 ecosystem to validate the full software stack against the FPGA fabric. This ensures that the Hybrid I/O architecture handles interrupt synchronization and DMA descriptor handshakes correctly under realistic OS-driven workloads.

\subsection{Results and Coverage Closure}
The accelerator achieved a 100\% pass rate over the full regression suite (25 suites, 8,915 test cases). Advanced verification results show:
\begin{itemize}
    \item \textbf{No Protocol Violations:} Zero SVA failures detected across all stress testing cycles, including AXI write-back verification.
    \item \textbf{ISA Coverage:} 21/21 operations verified including RELU and MAX extensions for ANN/LPR.
    \item \textbf{Functional Coverage:} 100\% closure on ISA opcode combinations, NoC port-to-port routing patterns, and DMA bi-directional data paths.
\end{itemize}
The robust verification environment confirms the IP is ready for physical deployment on the Zynq-7000 XC7Z020, and the License Plate Recognition demo has been validated end-to-end on a Haoyue Zynq-7020 board running PetaLinux 2022.2.
